\subsection{CUDA Implementation}
\label{sec:cuda_implementation}

The supplied ASR pipeline provides CUDA implementations for rotary
positional embeddings, reductions, general matrix multiplication
(GEMM), attention, and auxiliary operations. We completed the
required kernels and investigated the performance of the GEMM
implementation in particular.

\paragraph{Tiled matrix multiplication.}
For an input matrix $X\in\mathbb{R}^{M\times K}$ and weights
$W\in\mathbb{R}^{N\times K}$, a linear layer computes
\begin{equation}
Y_{m,n} =
f\left(\sum_{k=0}^{K-1}X_{m,k}W_{n,k}+b_n\right),
\end{equation}
where $f$ is the identity or GELU activation. The weight layout
matches PyTorch's convention of one row per output feature.

Our GEMM kernel assigns one output element to each thread in a
$T\times T$ thread block, with $T=16$ in the final implementation.
Each block processes the reduction dimension in tiles of width $T$.
Input and weight tiles are loaded into shared memory and reused
across the block. The weight tile uses a padded shared-memory
layout of $T\times(T+1)$ elements to mitigate bank conflicts.

For $M=N=K=1024$, a naive implementation performing independent
dot products requests $2MNK=2{,}147{,}483{,}648$ operand loads.
With ideal tile reuse, the tiled implementation requests approximately
$2MNK/T=134{,}217{,}728$ loads for $T=16$, a theoretical
sixteenfold reduction. This counts load requests rather than
measured device-memory traffic; caching and hardware effects
can change actual memory transfers.

Out-of-bounds tile elements are replaced with zeros. All threads,
including those without a valid output element, participate in
both synchronization barriers: the first ensures that shared-memory
loads finish before computation, and the second prevents the next
tile from overwriting values still in use.

\paragraph{Fused epilogue.}
The accumulated output remains in a register while optional bias
addition and GELU are applied, avoiding an additional activation
kernel. For encoder \texttt{fc1}, with output dimensions
$1500\times5120$, eliminating a separate FP32 read and write
saves $2(1500)(5120)(4)=61{,}440{,}000$ bytes of global-memory
traffic, excluding cache effects and launch overhead.


\paragraph{Attention and online softmax.}
Our attention kernel processes 32 keys per tile,
with four warps per block and one query per warp.
Keys and values are staged in shared memory while
each query maintains a running maximum, normaliser,
and weighted output accumulator. This follows the
online-softmax principle underlying FlashAttention
\citep{dao2022flashattention}, avoiding materialisation
of the full attention-score matrix. The kernel supports
grouped-query attention by mapping query heads to
shared key/value heads and applies causal masking
where required. Unlike production FlashAttention
kernels, it uses per-lane scalar dot products rather
than optimised matrix operations, limiting throughput
despite reduced intermediate storage.

\paragraph{Kernel selection.}
The tiled kernel is compared against a naive implementation and
NVIDIA cuBLAS. cuBLAS can exploit more advanced tiling, instruction
selection, and hardware-specific optimisations than this
straightforward shared-memory kernel. We additionally compare
tile sizes $T\in\{8,16,32\}$ and evaluate end-to-end configurations
that selectively replace custom CUDA operations with PyTorch
implementations.
